\begin{tikzpicture}[
    node distance=5mm and 6mm,
    block/.style={draw, rounded corners=2pt, align=center, minimum height=9mm, inner sep=4pt, font=\small},
    group/.style={draw, rounded corners=3pt, inner sep=5pt},
    flow/.style={-{Latex[length=2.2mm]}, semithick}
]
\node[block, fill=gray!10, minimum width=2.9cm] (ingress) {业务请求\\文本/长上下文/多模态/工具调用};
\node[block, fill=blue!8, minimum width=2.7cm, right=of ingress] (service) {服务入口与接口层\\OpenAI 兼容/鉴权/监控};
\node[block, fill=blue!8, minimum width=2.7cm, right=of service] (scheduler) {调度与运行时层\\batch 组织/SLO/bucket/回退};
\node[block, fill=green!8, minimum width=2.4cm, below=of scheduler] (cache) {KV Cache 与状态层\\前缀复用/迁移/驱逐};
\node[block, fill=green!8, minimum width=2.4cm, left=of cache] (execute) {执行层\\prefill/decode/并行策略};
\node[block, fill=green!8, minimum width=2.4cm, right=of cache] (backend) {后端接口层\\图模式/算子/能力门控};
\node[group, fit=(execute)(cache)(backend), label={[font=\small]above:运行时核心}] (runtime) {};
\node[block, fill=orange!10, minimum width=2.6cm, right=18mm of backend] (compiler) {编译器与运行时\\CANN/MindSpore/插件栈};
\node[block, fill=orange!10, minimum width=2.3cm, below=of compiler] (device) {国产硬件平台\\NPU/GPU/驱动/通信};

\draw[flow] (ingress) -- (service);
\draw[flow] (service) -- (scheduler);
\draw[flow] (scheduler) -- (backend);
\draw[flow] (scheduler) -- (execute);
\draw[flow] (scheduler) -- (cache);
\draw[flow] (execute) -- (cache);
\draw[flow] (cache) -- (backend);
\draw[flow] (backend) -- (compiler);
\draw[flow] (compiler) -- (device);
\draw[flow] (device.west) |- (backend.east);
\draw[flow, dashed] (cache.west) |- (service.south);

\node[font=\scriptsize, align=left, below=2mm of ingress] {真实压力来自多阶段请求异质性};
\node[font=\scriptsize, align=left, below=2mm of compiler] {平台差异集中暴露在编译、算子与回退边界};
\end{tikzpicture}